\section{Introduction}\label{sec:intro}

\IEEEPARstart{S}{atellite}-based free-space optical (FSO) channels are the most credible medium for intercontinental quantum key distribution (QKD): they avoid the exponential photon-loss bound of optical fiber and exploit the near-vacuum of space for low excess noise \cite{liao2017micius,bedington2017progress,sidhu2021advances,pirandola2020advances}. The 2017 demonstration of decoy-state BB84 between the Micius satellite and ground stations in China \cite{liao2017micius}, the simultaneous entanglement distribution at $1{,}200$ km baseline \cite{yin2017satellite}, and the relayed intercontinental key exchange of \cite{liao2018satellite} together establish satellite QKD as a near-term technology rather than a research curiosity. Recent progress is summarized in \cite{sidhu2021advances,pirandola2020advances}. Two reviews frame the present work in particular: \cite{xu2020rmp} surveys how device imperfections, rather than protocol design, set the security of practical QKD, which is the regime our parameter-control problem lives in; and \cite{lu2022micius} reviews the Micius programme, still the only satellite platform on which these protocols have been demonstrated end to end, and the source of the link budgets that make a multi-user relay architecture worth optimizing at all.

Among the QKD families adapted to FSO, the subcarrier-intensity-modulation BPSK with dual-threshold direct detection (SIM-BPSK/DT-DD) scheme proposed by Vu, Pham, Dang, Le, and Pham \cite{vu2020access,vu2022access,vu2023photonics} is attractive because it relies on entirely classical components (an EDFA-amplified GEO source, PIN photodetectors, and dual-threshold post-processing), yet is reported in \cite{vu2020access,vu2022access,vu2023photonics} to admit a positive secret-key rate against unauthorized-receiver attacks (URA) and beam-splitting attacks (BSA) \cite{lutkenhaus1999estimates,huttner1995quantum}, provided that the modulation depth $\mu$, the DT scale coefficients $\beta_A, \beta_i$, and the multi-user fading-randomness exclusion ratio $\chi$ are appropriately chosen. We state the scope of that guarantee precisely, because the present work inherits it rather than re-deriving it: the secret-key rate used throughout this paper is the asymptotic Csisz\'ar--K\"orner fraction $\max\{0, H_2(P_e) - H_2(\mathrm{QBER})\}$ per sifted bit, evaluated against the two attack models above. It carries no finite-key correction, no composable-security statement, and no claim of optimality against general coherent attacks; it is a sufficient-condition surrogate that we optimize, not a security proof that we establish. SIM-BPSK is well established for classical FSO links \cite{popoola2009bpsk,kaushal2017fso}, and DT-DD inherits its noise-budget structure from the standard outage-capacity analysis \cite{farid2007outage,sandalidis2008ber}.

The three papers in the lineage establish, in increasing generality, the relay-assisted single-link case \cite{vu2020access}, the entanglement-based two-user case \cite{vu2022access} (which adapts the BBM92 protocol \cite{bbm92} to the DT/DD measurement basis), and the GEO source serving $N$ ground users via dual LEO relays with a closed-form inclusion-exclusion term \cite{vu2023photonics}. In all three, the parameters $(\mu, \beta, \chi)$ are picked \emph{statically}: a single operating point is selected by reading the analytic performance plots at a representative geometry, and that point is then held throughout operation.

We show in this paper that for the realistic case of a LEO relay traversing its pass while serving an asymmetric cluster of ground users, the static-parameter assumption is the dominant performance bottleneck. Peak elevation, slant range, log-amplitude turbulence variance \cite{andrews2005laser,vetelino2007fade}, and the URA leakage geometry vary by orders of magnitude across a single 5 to 10 minute pass, and the optimal $\bp(t) = (\mu, \beta_A, \{\beta_i\}, \chi)$ drifts continuously with the time-varying state $\bs(t)$. Under realistic constellations propagated from real Two-Line Elements (TLE) via the SGP4 propagator \cite{hoots1980sgp4,vallado2006sgp4} on Starlink-density shells \cite{del2019leosat,leyva2021inter}, we further show that \emph{no} single static $\bp$ produces any secure pair-step over a typical $90$-minute observation window, so adaptation is a precondition for the system to operate at all rather than a marginal improvement.

The modulation-signalling branch of QKD that this paper belongs to has moved on since the three papers above, and in directions that sharpen what is and is not new here. Subcarrier-wave QKD has been analysed numerically for free-space ship-to-port links \cite{goncharov2025scw} and routed through a deployed metropolitan transport network \cite{tarabrina2023scw}, which establishes the signalling format outside the laboratory but in both cases for a \emph{single} link rather than a shared relay. On the transmitter side, direct intensity and phase modulation has been demonstrated as a QKD source \cite{liu2026intensity}, which is the hardware assumption our modulation depth parameterizes. The free-space continuous-variable line is surveyed in \cite{kargina2026cvqkd}; it shares our reliance on classical coherent components but replaces the dual-threshold decision with homodyne detection, so the parameter that dominates our problem, a detection threshold, has no counterpart there. What none of these address is the coupling we study: a single global modulation depth that every user in a cluster must share.

Learning-based parameter control for QKD is itself an active line, and the closest neighbour to this work is \cite{mohamed2026aml}, which couples a temporal convolutional forecaster with a proximal-policy-optimization agent to retune decoy intensities online for BB84, E91 and COW while keeping a composable-security constraint. The contrast is instructive in three ways. First, the physical layer differs: that work tunes photon-counting discrete-variable protocols, whereas the parameters here are a modulation depth and a pair of detection thresholds in an intensity-modulated direct-detection link, so the two parameter spaces do not overlap and a numerical head-to-head comparison would compare protocols rather than controllers. Second, that work is single-link; the coupling that makes the present problem hard, a single global modulation depth shared by every Alice-Bob pair in a cluster, does not arise there. Third, it learns online by reinforcement while we learn offline by supervision from a decomposed oracle, which buys determinism and an inspectable target at the cost of needing that oracle. A reinforcement-learning controller for the multi-user FSO setting is a worthwhile comparison that we leave to future work; building a fair one is a study in itself, and a weak re-implementation would be worse evidence than none.

\subsection{Contributions}

This paper makes the following contributions:

\begin{enumerate}
\item \textbf{Full multi-user joint problem.} We restate the GEO/multi-LEO/$N$-user system of \cite{vu2023photonics} together with the pointing-error/Gaussian-beam model of \cite{vu2020access} (dropped in \cite{vu2023photonics}) and derive the corresponding joint problem \textbf{(P)} over $\bp=(\mu, \beta_A, \{\beta_i\}, \chi)$ with URA, BSA, and an inter-user secrecy constraint. The problem is nonconvex with three sources of coupling: $\mu$ shared globally, $\chi$ entering every pair through a product field $\phi_i = \prod_{j\neq i}(1-P_c^j)$, and $\beta_A$ shared between Alice and every Bob. The inclusion--exclusion term of \cite[Eq.~22]{vu2023photonics} is generalized to asymmetric clusters in closed form.
\item \textbf{Principled two-level decomposition.} We expose the variable roles to obtain a two-level decomposition whose outer/inner split is exact (a relaxation-free reformulation of \textbf{(P)}): an outer two-dimensional search over $(\mu,\chi)$ wraps an inner mean-field block-coordinate solver over the thresholds. The symmetric-cluster fixed point provably recovers \cite[Eq.~24]{vu2023photonics}, and the decomposition matches a brute-force joint grid optimum on a two-user cluster to four decimals.
\item \textbf{Two-stage controller with KAN.} The decomposed oracle generates supervised pairs $(s(t)\!\to\!p^\star(t))$ for a two-stage controller: a global head consumes cluster-aggregate features and predicts $(\mu, \chi, \beta_A)$; a per-user shared head consumes $(s_i, \beta_A, \mu, \chi)$ and predicts $\beta_i$. We compare a Kolmogorov--Arnold Network (KAN) \cite{liu2024kan,liu2024kan2} against a regularized multi-layer perceptron and a linear residual baseline on a five-axis trade-off (regression accuracy, closed-loop secret-key retention, parameter count, inference latency, and interpretability). Deep learning for the classical physical layer \cite{oshea2017introduction,wang2017deep,ye2018deep} and reinforcement-learning control \cite{luong2019applications} are well studied; KAN is attractive here because its spline activations admit symbolic extraction \cite{liu2024kan2,cranmer2020discovering}, which yields a closed-form rule for $\beta^\star$.
\item \textbf{TLE-driven realistic-orbit validation.} An SGP4 path mirrors the analytic-orbit API so every script accepts \texttt{--tle} as a drop-in switch, with a synthetic Walker shell reproducing Starlink coverage statistics offline and a Celestrak bundle loadable via \texttt{--tle PATH}. Over a Hanoi ground station, the static baseline collapses to zero secure pair-steps on the Walker constellation, while the adaptive controller recovers a non-trivial fraction near overhead.
\item \textbf{Honest multi-axis evaluation.} We report results with explicit mean$\pm$std over five seeds, fair (regularized) MLP baselines, and a model-agnostic $\beta$ safety-margin recipe that lifts feasibility by 1.1--1.5$\times$ for every controller. KAN does not dominate every axis; its value is parameter efficiency plus an interpretable closed-form rule for $\beta$ with symbolic $R^2 \approx 0.97$. We disclose where the symbolic extraction fails (the $\mu^\star$ saturation regime), where the controllers transfer (single-link) and where they do not (multi-user cluster, the residual-learning rescue), and how the design generalizes to real-orbit operation.
\end{enumerate}
